\section{Coordinate densities}\label{app:geometry}

We prove the density bounds used to choose anchors and to turn coordinate
approximation into probability approximation. The first bounds follow by
fixing the incomparable coordinates and comparing sections of the remaining
convex region. Let $V$ be the set of coordinates, let $N=|V|$, and let
$K\subseteq\R^V$ be their feasible region.
For a coordinate $x$ and a real number $t$, write
\[
K_x(t)=\{(z_i)_{i\ne x}:z\in K,\ z_x=t\},
\qquad S_x(t)=\vol_{N-1}K_x(t).
\]
Here $K_x(t)$ is the $(N-1)$-dimensional section obtained by fixing
coordinate $x$ at $t$. The uniform distribution on $K$
therefore gives $x$ the density
\[
f_x(t)=\frac{S_x(t)}{\vol_N K}.
\]
These formulas define $f_x$ pointwise, and we use this definition throughout.
If some coordinates have already been fixed, we use the same
definitions in the remaining free coordinates. In dimension zero, the
volume of the singleton consisting of the empty assignment is one.

\begin{proof}[Proof of Lemma~\ref{lem:fiber}]
We may first fix any prescribed set of coordinates incomparable with $x$,
at values leaving a positive-volume slice; write $f_x$ for its coordinate
density. Fixing a coordinate turns its incident inequalities
into interval bounds. These bounds propagate along directed paths, so
fixing an incomparable coordinate leaves $x$'s effective endpoints
unchanged. By Lemma~\ref{lem:interior}, its projection interval remains
$(\ell,r)$. The same argument applies successively as more incomparable
coordinates are fixed.

Now fix all the remaining coordinates incomparable with $x$. For each
choice leaving a nonempty slice, the remaining feasible region $C$ is
open and convex, has dimension $d\le B+1$, and projects onto $(\ell,r)$
in coordinate $x$. Write $S(t)$ for its section volume at $x=t$ and
$g(t)=S(t)/\vol C$ for the resulting density.

Choose points in $\overline C$ whose $x$-coordinates are $\ell$ and $r$.
Such points exist because $C$ is bounded. For $\ell<t<s<r$, join every
point of the section at $t$ to the point at $r$. Since $C$ is open and
convex, the image at $x=s$ lies in the section of $C$ there. Each of its
$d-1$ directions is scaled by $(r-s)/(r-t)$, so
\begin{equation}\label{eq:push-upper}
S(s)\ge\left(\frac{r-s}{r-t}\right)^{d-1}S(t)
\qquad(t<s).
\end{equation}
Joining instead to the point at $\ell$ gives
\begin{equation}\label{eq:push-lower}
S(s)\ge\left(\frac{s-\ell}{t-\ell}\right)^{d-1}S(t)
\qquad(s<t).
\end{equation}
Normalize these inequalities and integrate over $s$ on both sides of $t$:
\[
1\ge g(t)\left(
\int_\ell^t\left(\frac{s-\ell}{t-\ell}\right)^{d-1}ds+
\int_t^r\left(\frac{r-s}{r-t}\right)^{d-1}ds\right)
=\frac{Wg(t)}{d}.
\]
Thus $g(t)\le d/W\le(B+1)/W$.

For the lower bound, choose $t_0$ with $g(t_0)\ge1/W$, which exists
because $g$ integrates to one over an interval of length $W$. For any
$s\in[\ell+W/4,r-W/4]$, the appropriate ratio in
\eqref{eq:push-upper} or \eqref{eq:push-lower} is at least $1/4$.
Consequently
\[
g(s)\ge4^{-(d-1)}g(t_0)\ge\frac{4^{-B}}{W}.
\]
The case $s=t_0$ satisfies the same bound directly.

Finally, average over the additional incomparable coordinates fixed in
the second paragraph. If $\omega$ denotes their values, $C_\omega$ the
remaining region, and $g_\omega$ its coordinate density, integration of
each section gives the pointwise identity
\[
f_x(t)=
\frac{\int \vol(C_\omega)g_\omega(t)\,d\omega}
     {\int \vol(C_\omega)\,d\omega}.
\]
Empty slices contribute zero. The bounds on $g_\omega$ are uniform in
$\omega$, so they hold for $f_x$ at every $t$. This proves both the
unconditional bounds and the bounds after the initial conditioning.
\end{proof}

The coordinate densities above are log-concave pointwise: the uniform
density on $K$ is log-concave, and integrating out the other coordinates
preserves this property by Pr\'ekopa's marginal theorem
\cite[Theorem~8]{prekopa}. We finish with the density bound used for
positions having substantial probability on each side.

\begin{proof}[Proof of Lemma~\ref{lem:central-quantiles}]
For $s\le c\le t$, log-concavity gives
\[
f(s)f(t)\le f(c)f(s+t-c).
\]
Indeed, $c$ and $s+t-c$ are complementary convex combinations of $s,t$;
multiply the two log-concavity inequalities. Fixing $t\ge c$ and
integrating over $s\le c$ gives
\[
p f(t)\le f(c)\int_{-\infty}^c f(s+t-c)\,ds\le f(c).
\]
Here the first inequality uses the mass at least $p$ to the left of $c$,
and the last uses that a density integrates to one. Fixing $s\le c$ and
integrating over $t\ge c$ instead uses the right-hand mass and gives
$p f(s)\le f(c)$. Taking the supremum over both sides of $c$ proves the
claim.
\end{proof}
